\subsection{Bounded subsets of a topological vector space}

DEFINITION 3. --- Let E be a topological vector space. A subset A of E is said to be bounded if it is absorbed by every neighbourhood of 0 in E (I, p.~7, def. 4).

In order that A be bounded, it is sufficient that A be absorbed by every neighbourhood of a fundamental system of neighbourhoods of 0. Since there exists a fundamental system of balanced neighbourhoods of 0 (I, p.~7, prop. 4), this is the same as saying that, for every neighbourhood V of 0 in E, there exists \(\lambda \in K\) such that \(A \subset \lambda V\) .

Suppose the topology of E is defined by a fundamental system \(\Gamma\) of semi-norms (II, p.~3); then a subset A of E is bounded if and only if every semi-norm \(p \in \Gamma\) is bounded on A.

In particular, if E is a semi-normed space, a subset A of E is bounded if and only if it is contained in a ball. In other words, if E is normed this means that A is bounded for the metric space structure of E (GT, IX, § 2, No.~2).

Remarks. --- 1) If E is a semi-normed space, the balls form a fundamental system of bounded neighbourhoods of 0 in E. Conversely, if E is a locally convex topological vector space, and if there exists a bounded neighbourhood of 0 in E, this neighbourhood contains a convex balanced neighbourhood W, and the gauge of W is then a semi-norm defining the topology of E.

\pandocbounded{\includegraphics[keepaspectratio]{images/part001_149bf36b.jpg}}

Thus, if E is locally convex and metrizable, and if its topology cannot be defined by a single norm, then there exists no distance on E defining its topology and such that the bounded subsets for \(d\) (GT, IX, § 2, No.~2) are the bounded subsets of E. More precisely, for every distance \(d\) on E, which is translation invariant and which defines the topology of E, the bounded subsets of E are bounded for \(d\) (III, p.~38, exerc. 3), but the converse is false.

\begin{enumerate}
\def\labelenumi{\arabic{enumi})}
\setcounter{enumi}{1}
\item
  Let M be a vector subspace of E endowed with the induced topology. In order that a subset of M be bounded in M, it is necessary and sufficient that it be bounded in E.
\item
  Let N be the intersection of all neighbourhoods of 0 in E, so that \(\tilde{E} = E/N\) is the Hausdorff vector space associated with E. Then N is bounded; if \(\pi: E \to \tilde{E}\) is the canonical homomorphism then a subset B of E is bounded if and only if \(\pi(B)\) is bounded.
\item
  Let E be a Hausdorff locally convex space; then for every \(x \neq 0\) in E, there exists a continuous semi-norm p such that \(p(x) \neq 0\) ; this semi-norm is not bounded on the real half-line \(R_{+}.x\) generated by x. Hence no non-null subspace of E is bounded. In particular, a bounded subset does not contain any line.
\end{enumerate}

DEFINITION 4. --- Let E be a locally convex space. A bornology B on E is said to be adapted to E, if it is convex, is composed of bounded subsets of E and if the closure of every set of B belongs to B.

PROPOSITION 1. --- Let E be a locally convex space. The set of bounded subsets of E is an adapted bornology.

We need to establish the following properties :

\begin{enumerate}
\def\labelenumi{\alph{enumi})}
\item
  If \(\mathbf{B}\) is a bounded subset of \(\mathbf{E}\) , every subset of \(\mathbf{B}\) is bounded.
\item
  The union of two bounded subsets is bounded.
\item
  Every set that is homothetic to a bounded set is bounded.
\item
  The closed convex balanced envelope (II, p.~13) of a bounded subset is bounded.
\end{enumerate}

If p is a continuous semi-norm on E, the balls of p are convex, balanced, closed and the set homothetic to a ball is a ball. Hence, if p is bounded on two subsets X and Y of E, it is also bounded on the closed convex balanced envelope of \(X \cup Y\) , and on the sets homothetic to these. This establishes properties b), c) and d), and a) is obvious.

DEFINITION 5. --- Let E be a locally convex space. The set of all bounded subsets of E is called the canonical bornology of E.

If B is a set of bounded subsets of E, then there exists a smallest bornology \(\tilde{B}\) adapted to E and containing B. The sets of \(\tilde{B}\) are those that are contained in a set homothetic to the closed convex balanced envelope of a finite union of sets of B.

Every adapted bornology is contained in the canonical bornology.

PROPOSITION 2. --- In a locally convex space E, every precompact set is bounded. Let A be a precompact subset of E, and V be a convex balanced neighbourhood of 0. There exists a finite sequence \((a_{i})_{1\leqslant i\leqslant n}\) of points of A such that

\[
\mathrm{A} \subset \bigcup_ {1 \leqslant i \leqslant n} (a _ {i} + \mathrm{V}).
\]

Since \(B = \{a_{1}, ..., a_{n}\}\) is bounded, there exists a scalar \(\lambda\) such that \(0 < \lambda < 1\) and \(\lambda B \subset V\) ; we have \(\lambda A \subset \lambda B + \lambda V \subset V + V\) , from which the proposition follows.

COROLLARY. --- In a locally convex space, the set of points of a Cauchy sequence is bounded.

In fact, this set is precompact (GT, II, § 4, No.~2).

Remark 5. --- In general the bounded subsets of a locally convex space E are not all precompact (for example, if E is an infinite dimensional normed space, its unit ball is not compact (I, p.~15, th. 3)). However, this is so if E is a weak space (II, p.~42) : for the Hausdorff topological vector space associated with E is then isomorphic to a subspace of a product \(K^{I}\) whose bounded subsets are precompact (cf.~III, p.~4, cor. 2).

For other examples, see IV, p.~18.

PROPOSITION 3. --- Let A be a subset of a locally convex space E. Suppose that A is bounded; then for every sequence \((x_{n})\) of points of A and for every sequence \((\lambda_{n})\) of scalars tending to 0, the sequence \((\lambda_{n}x_{n})\) tends to 0. Conversely, if there exists a sequence \((\lambda_{n})\) of non-zero scalars such that for every sequence \((x_{n})\) of points of A, the sequence \((\lambda_{n}x_{n})\) is bounded, then A is bounded.

Suppose that A is bounded. If \((\lambda_{n})\) is a sequence of scalars tending to 0, and V is a neighbourhood of 0, we have \(\lambda_{n}A \subset V\) whenever n is large enough, and the first assertion follows.

Conversely, if A is not bounded and if \((\lambda_{n})\) is a sequence of scalars \(\neq 0\) , then there exists a continuous semi-norm p and a sequence \((x_{n})\) of points of A, such that \(p(x_{n}) \geqslant \frac{n}{|\lambda_{n}|}\) . We have then that \(p(\lambda_{n}x_{n}) \geqslant n\) , and the sequence \((\lambda_{n}x_{n})\) is not bounded.

COROLLARY. --- A subset A of E is bounded if and only if every countable subset of A is bounded.

